\documentclass[11pt]{article}
\usepackage[a4paper,margin=21mm]{geometry}
\usepackage[no-math]{fontspec}
\setmainfont{Latin Modern Roman}
\newfontfamily\CJKfont{Noto Serif CJK SC}
\XeTeXlinebreaklocale "zh"
\XeTeXlinebreakskip=0pt plus 1pt
\usepackage{amsmath,amssymb,amsthm,amscd,graphicx,color}
\usepackage{xurl}
\usepackage[unicode,hidelinks]{hyperref}
\allowdisplaybreaks
\setlength\emergencystretch{3em}
\numberwithin{equation}{section}

\font\teneur=eurm10\font\seveneur=eurm7
\font\teneus=eusm10\font\seveneus=eusm7
\font\sans=cmss10

\renewcommand{\d}{\mathrm{d}}%exterior derivative
\newcommand{\cL}{{\mathcal L}}
\newcommand{\cI}{{\mathcal I}}
\newcommand{\bbR}{{\mathbb R}}
\newcommand{\bbC}{{\mathbb C}}
\newcommand{\bbP}{{\mathbb P}}
\newcommand{\E}{{\mathrm e}}
\newcommand{\Gr}{\operatorname{Gr}}
\newcommand{\Leg}{\operatorname{Leg}}
\newcommand{\Lag}{\operatorname{Lag}}
\newcommand{\Ss}{\mathsf{S}}
\newcommand{\w}{\wedge}

\numberwithin{equation}{section}

\newtheorem{Theorem}{Theorem}[section]
\newtheorem*{Theorem*}{Theorem}
\newtheorem{Corollary}[Theorem]{Corollary}
\newtheorem{Lemma}[Theorem]{Lemma}
\newtheorem{Proposition}[Theorem]{Proposition}
 { \theoremstyle{definition}
\newtheorem{Definition}[Theorem]{Definition}
\newtheorem{Note}[Theorem]{Note}
\newtheorem{Example}[Theorem]{Example}
\newtheorem{Remark}[Theorem]{Remark} }

\begin{document}
\begin{center}\Large Strictly convex Pfaff-Darboux representations adapted to a positive Legendrian foliation for a torsion-free affine connection\end{center}
\noindent\textbf{Stable ID:} 1149\quad\textbf{Author:} Robert L. Bryant\par
\noindent\textbf{Work:} On the Convex Pfaff-Darboux Theorem of Ekeland and Nirenberg\par
\noindent\textbf{Source:} \url{https://sigma-journal.com/2023/060/}\par
\noindent\textbf{Version:} Official SIGMA 19 (2023) article 060, DOI 10.3842/SIGMA.2023.060; journal PDF and native TeX acquired 2026-09-30, exact bytes pinned by SHA256\par
\noindent\textbf{Rights:} CC BY 4.0. Authors retain copyright. \url{https://creativecommons.org/licenses/by/4.0/}. Reuse and typesetting adaptation; original language and mathematical content preserved.\par
\noindent\textbf{Difficulty (prior selection preserved):} {\CJKfont H2: The proof upgrades positivity only along the prescribed Legendrian leaves to strict positivity of every coordinate Hessian on the ambient tangent space. It successively changes the Pfaff coordinates by quadratic corrections, a one-variable convex reparametrization and controlled positive linear mixtures. The final perturbation makes all coefficient functions positive while preserving the given foliation and the exact one-form representation, an additional construction beyond the ordinary Pfaff-Darboux normal form.}\par
\medskip\noindent\textbf{Editorial boundary note (not author text):} Complete original principal Theorem3.3 and the entire foliation-adapted convex/positive coordinate construction. All fixed Pfaff-rank/nonvanishing assumptions, canonical kernel/characteristic subbundles, Legendrian leaves and classical adapted normal-form conventions, torsion-free affine/Hessian and symmetric derivative definitions, positive leaf restriction and necessary-condition proof are retained. The full new proof includes rescaling of the one-form, quadratic correction along the hyperplane, one-variable strictly convex reparametrization, positive linear mixtures and the final small-epsilon shift making every coefficient positive while preserving the foliation. The classical Pfaff-Darboux theorem is precisely stated prior background; no recursive reproof or replacement is added. One expressly announced strengthened/generalized main outcome; the unadapted corollary and Euclidean application are unselected, with the necessary-condition proposition/context receiving no extra counts. Literal source superscript/subscript y notation and all coefficients/signs are preserved.\par
\medskip\hrule\medskip\noindent\textbf{Author text begins below.}\par
\setcounter{section}{1}
% Original source lines 175--315
\section{Classical Pfaff--Darboux theorems}
Let $\omega$ be a smooth $1$-form on an $n$-manifold~$M$
that, for some integer $k>0$, satisfies
\begin{equation}\label{eq: omegaPfaff_k}
\text{$\omega\w(\d\omega)^k$ vanishes identically on $M$}
\end{equation}
while
\begin{equation}\label{eq: omega_nondegenerate}
\text{$\omega\w(\d\omega)^{k-1}$ is nowhere vanishing on $M$.}
\end{equation}
The integer~$k{-}1$ is known as the \emph{Pfaff rank}
of~$\omega$~\cite[Chapter II, Section~3]{MR1083148}.
Note that $k\le \tfrac12(n{+}1)$. When $k=\tfrac12(n{+}1)$,
$\omega$ is said to be a \emph{contact form} on~$M$.

When $\omega$ satisfies~\eqref{eq: omegaPfaff_k}
and~\eqref{eq: omega_nondegenerate}, so does $\tilde\omega=f\omega$
for any nonvanishing function~$f$ on~$M$,
since $\tilde\omega\w(\d\tilde\omega)^{r-1} = f^r \omega\w(\d\omega)^{r-1}$
for all integers~$r>0$.

\subsection{Canonical subbundles}
An $\omega$ satisfying~\eqref{eq: omegaPfaff_k}
and~\eqref{eq: omega_nondegenerate} defines a \emph{kernel subbundle}~$K
= \omega^{-1}(0)\subset TM$ of corank~$1$
and a subbundle $A\subset K$ of corank $2(k{-}1)$ in~$K$
by the rule that, for each $m\in M$,
\begin{equation*}
A_m = \{ v\in K_m\mid \mathrm{d}\omega(v,w) = 0,\, \forall w\in K_m \}.
\end{equation*}
Replacing~$\omega$ by $\tilde\omega= f \omega$ for any nonvanishing
function~$f$ does not change~$K$ or $A$.

If $k>1$, then $K\subset TM$ is not an integrable plane field,
but the subbundle $A\subset TM$ is always integrable,
since it is the Cauchy characteristic plane field
of the differential ideal~$\cI$ generated by~$\omega$
(see~\cite[Chapter II, Proposition~2.1]{MR1083148}).
In the contact case, i.e., when $n=2k{-}1$ (which is,
in some sense, generic), one has~$A=(0)$.

There is a nondegenerate, skew-symmetric bilinear pairing
$B_\omega \colon K{/}A\times K{/}A\to\bbR$
defined by
\[
B_\omega\bigl(v{+}A_m,w{+}A_m\bigr) = \mathrm{d}\omega(v,w),
\]
when $v,w\in K_m$. It satisfies $B_{f\omega} = fB_\omega$
for any nonvanishing function~$f$ on~$M$.

Note that any subspace~$W\subset T_m$ on which both $\omega$ and $\d\omega$
vanish, must, first of all, satisfy $W\subset K_m$
(since $\omega$ vanishes on~$W$), and, second, must have codimension
at least $k{-}1$ in $K_m$, since $\d\omega$, as a skew-symmetric form
on $K_m$, has Pfaff rank~$k{-}1$.
Moreover, if $W$ does have codimension $k{-}1$ in $K_m$,
then it must contain $A_m$, so that $W/A_m$ is a null subspace of $B_\omega$.

\subsection{Legendrian submanifolds and Grassmannians}
Any submanifold~$L\subset M$ that satisfies~$L^*\omega=0$,
i.e., an \emph{integral manifold of~$\omega$},
must also satisfy $L^*\d\omega=0$
and hence, by the above linear algebra discussion,
must have codimension at least~$k$ in~$M$.
When $L\subset M$ is an integral manifold of~$\omega$ of codimension~$k$,
it is said to be an \emph{$\omega$-Legendrian submanifold}.

In particular, if $L\subset M$ is $\omega$-Legendrian,
then, for each $m\in L$, the tangent space $T_mL$
satisfies $A_m\subset T_mL\subset K_m$
while $B_\omega$ vanishes identically on $T_mL/A_m\subset K_m/A_m$.

This motivates defining
the \emph{Legendrian Grassmannian}~$\Leg_m(\omega)\subset\Gr^k(T_mM)$
to be the set of subspaces~$W\subset K_m$
that have codimension~$k$ in $T_mM$
and on which both~$\omega$ and~$\d\omega$ vanish.
By the above remarks, it follows that $\Leg_m(\omega)$
can be canonically identified with the
\emph{Lagrangian Grassmannian}~$\Lag(K_m/A_m)\subset\Gr^{k-1}(K_m/A_m)$
consisting of the $(k{-}1)$-dimensional subspaces of~$K_m/A_m$
on which $B_\omega$ vanishes.
Hence, $\Leg_m(\omega)$ is naturally a smooth manifold
of dimension~$\tfrac12k(k{-}1)$.
Moreover, the (disjoint) union
\[
\Leg(\omega) = \bigcup_{m\in M} \Leg_m(\omega)\subset \Gr^k(TM)
\]
is a smooth subbundle,
and $\Leg(f\omega)=\Leg(\omega)$ for all nonvanishing~$f$.

\subsection{A local normal form}
One version of the \emph{Pfaff--Darboux theorem}
\cite[Chapter II, Theorem~3.1]{MR1083148} states that,
when ${\omega\in\Omega^1(M)}$ satisfies~\eqref{eq: omegaPfaff_k}
and~\eqref{eq: omega_nondegenerate},
each $m\in M$
has an open neighborhood~$U\subset M$
on which there exist smooth functions~$u^1,\ldots,u^k$
and $a_1,\ldots,a_k$ (with not all $a_i$ simultaneously vanishing)
so that
\begin{equation}\label{eq: PfaffDarboux_normalform}
U^*\omega = a_1\,\d u^1 + \cdots + a_k\,\d u^k .
\end{equation}
Moreover,
the mapping~$(u,[a])\colon U\to\bbR^k\times\mathbb{RP}^{k-1}$
is a submersion. (In fact, the kernel subbundle of the differential
of this mapping is the restriction of~$A$ to~$U$.)

Conversely, the existence of functions $u^i$ and $a_i$ for~$1\le i\le k$
on an open set~$U\subset M$ satisfying~\eqref{eq: PfaffDarboux_normalform}
with the $a_i$ not all simultaneously vanishing and having the
property that $(u,[a]) \colon U\to\bbR^k\times\mathbb{RP}^{k-1}$
be a submersion implies that both~\eqref{eq: omegaPfaff_k}
and~\eqref{eq: omega_nondegenerate} hold on~$U$.

\subsection{Geometry of the normal form}
It will be useful to have a geometric interpretation
of the Pfaff--Darboux theorem.
Now, in the representation~\eqref{eq: PfaffDarboux_normalform},
the functions~$u^i$ have independent differentials,
i.e., $\d u^1\w\cdots\w\d u^k$ does not vanish on~$U$.
Consequently, the simultaneous level sets of the functions~$u^i$
define a foliation~$\cL$ of~$U\subset M$
by $\omega$-Legendrian submanifolds, i.e., an $\omega$-Legendrian foliation.

Conversely, given an $\omega$-Legendrian foliation~$\cL$
on an open subset~$V\subset M$, each point~$m\in V$ will
have an open neighborhood~$U\subset V$ in which the leaves of~$\cL$
are the fibers of a submersion~${u = (u^i) \colon U\to\bbR^k}$.
Since~$\omega$ vanishes when pulled back to any fiber of~$u$,
it follows that there exists a mapping~$a = (a_i) \colon U\to\bbR^k$
such that $U^*\omega = a_1\,\d u^1 + \cdots + a_k\,\d u^k$.

Thus, a geometric interpretation of the Pfaff--Darboux theorem
is the statement that,
when $\omega\in\Omega^1(M)$ satisfies~\eqref{eq: omegaPfaff_k}
and~\eqref{eq: omega_nondegenerate}, each point~$m\in M$
has an open neighborhood~$U\subset M$
on which there exists an $\omega$-Legendrian foliation.



% Original source lines 342--532
\section{Convexity and affine manifolds}

\subsection{Classical convexity}
When~$M=\bbR^n$, there is a notion of \emph{strict convexity}
of a function~$u$, which is the condition
that the Hessian quadratic form~$H(u)$ be positive definite, where
\begin{equation}\label{eq: flat_Hessian}
H(u) = {\frac{\partial^2 u}{\partial x^i\partial x^j}}
\,\d x^i\otimes\d x^j
\end{equation}
and where $x^1,\ldots,x^n$
are the usual affine linear coordinates in~$\bbR^n$.
Note that strict convexity is an affine-invariant notion on~$\bbR^n$.

Motivated by applications in economics,
Ekeland and Nirenberg~\cite{MR2023129}
asked what further conditions one must impose on an~$\omega\in\Omega^1(\bbR^n)$
satisfying~\eqref{eq: omegaPfaff_k} and~\eqref{eq: omega_nondegenerate}
in order to know that one can choose the functions~$u^j$
and $a_j$ in the representation~\eqref{eq: PfaffDarboux_normalform}
so that the $u^j$ be strictly convex and the $a_j$ be positive.
It is not hard to show, by example, that \emph{some}
further condition on~$\omega$ is necessary
to guarantee the existence of such a convex representation.
(See the discussion at the beginning of Section~\ref{ssec: poscond}.)

They showed that two earlier articles~\cite{MR1655519,MR1652709}
claiming to provide such necessary and sufficient conditions were flawed
(indeed, they exhibited counterexamples to the claims of these articles)
and then produced their own condition,
which they showed to be necessary and sufficient.

In this note, I will show that their main result,
properly formulated, holds good on an $n$-manifold~$M$
endowed with a torsion-free affine connection,
not just on~$\bbR^n$ endowed with the (flat) affine connection
it inherits as a vector space.

\subsection{Affine connections and convexity}

Let~$\nabla$ be a torsion-free affine connection on an $n$-manifold~$M^n$,
i.e., $\nabla$ is a first-order, linear differential operator
\begin{equation*}
\nabla \colon \ \Omega^1(M)\to \Omega^1(M)\otimes \Omega^1(M)
\end{equation*}
that obeys the Leibnitz rule
\begin{equation*}
\nabla (f\eta) = \d f\otimes\eta + f \nabla(\eta)
\end{equation*}
for all smooth functions~$f$ on~$M$ and smooth $1$-forms~$\eta$ on~$M$.
The assumption that~$\nabla$ be torsion-free
is the condition that the associated (second-order)
\emph{Hessian operator}~$H(u)=\nabla(\d u)$
be a symmetric $(0,2)$-tensor for each smooth function~$u$ on~$M$.

A smooth function~$u$ on~$M$ is said to be {\it strictly $\nabla$-convex}
if, as a quadratic form,~$H(u)$ is positive definite at every point of~$M$.

When~$M=\bbR^n$ and~$\nabla$ is the standard (flat) connection,
satisfying~$\nabla(\d x^i) = 0$ for all of the coordinate functions~$x^i$,
then~$H(u)$ is the usual Hessian tensor~\eqref{eq: flat_Hessian},
and this notion of convexity is simply the classical one.

In the more general case, when $x = (x^i)\colon U\to \bbR^n$ is a local
coordinate chart, one has
\begin{equation*}
H\big(x^k\big) = \nabla\big(\d x^k\big) = \Gamma^k_{ij}\,\d x^i\otimes\d x^j,
\end{equation*}
where $\Gamma^k_{ij}=\Gamma^k_{ji}\in C^\infty(U)$
are the coefficients of the connection $\nabla$ relative
to the coordinate chart~$x=(x^i)$.
The general coordinate formula for~$H$ then becomes
\begin{equation*}
H(u) = \left(\frac{\partial^2 u}{\partial x^i\partial x^j}
+ \Gamma^k_{ij} \frac{\partial u}{\partial x^k}\right)
\,\d x^i\otimes\d x^j.
\end{equation*}
Thus, $\nabla$-convexity of $u$ is expressible in terms of a condition
on the $2$-jet of~$u$, slightly more general than the condition for
classical convexity.

Adopting the usual conventions
\begin{equation*}%\label{eq: altsym_decomp}
\alpha\w\beta
= \tfrac12(\alpha\otimes\beta-\beta\otimes\alpha), \qquad
\alpha\circ\beta
= \tfrac12(\alpha\otimes\beta+\beta\otimes\alpha),
\end{equation*}
one sees that, for a $1$-form~$\omega$ of the form
\begin{equation}\label{eq: P-Drepresentation}
\omega = a_1\,\d u^1 + \cdots + a_k\,\d u^k,
\end{equation}
one has (using the summation convention)
\begin{align*}
\nabla\omega
&= \d a_i\otimes\d u^i + a_i \nabla(u^i) = \d a_i\otimes\d u^i + a_i H(u^i)\\
&= \d a_i\w\d u^i + \d a_i\circ\d u^i + a_i H(u^i)\\
&= \d\omega + \Ss^\nabla\omega,
\end{align*}
where I have introduced
the notation~$\Ss^\nabla\omega$ to denote the {\it symmetrization\/}
of~$\nabla(\omega)$. Thus,~$\Ss^\nabla\omega
= \nabla\omega- \d\omega$
is a well-defined quadratic form on~$M$.
(Of course, the linear, first-order differential operator~$\Ss^\nabla$
depends on~$\nabla$.)


\subsection{A positivity condition}\label{ssec: poscond}
The quadratic form $S^\nabla\omega$ provides some insight
into the question of whether a $\nabla$-convex Pfaff--Darboux
representation of $\omega$ is possible.

\begin{Proposition}\label{prop: ness cond}
Suppose that $\omega\in\Omega^1(M)$ satisfies~\eqref{eq: omegaPfaff_k}
and~\eqref{eq: omega_nondegenerate}. If there exist positive
functions $a_i$ and $\nabla$-convex functions $u^i$ for $1\le i\le k$
such that~\eqref{eq: P-Drepresentation} holds, then $S^\nabla\omega$
is positive definite on the leaves of the foliation $\cL$ defined by
$\d u^1 = \d u^2 = \cdots = \d u^k = 0$.
\end{Proposition}

\begin{proof}
Since $\Ss^\nabla\omega =\d a_i\circ \d u^i + a_i H(u^i)$,
it follows that, when restricted to the plane field ${N\subset TM}$
defined by~$\d u^1 = \d u^2 = \cdots = \d u^k = 0$,
the terms~$\d a_i\circ \d u^i$ in~$\Ss^\nabla\omega$
vanish. Thus, $S^\nabla\omega = a_i H(u^i)$
as quadratic forms on $N$. By the positivity of the $a_i$
and the $\nabla$-convexity of the $u^i$, it follows that
$S^\nabla\omega$ is positive definite on~$N$.\footnote{It is worth pointing out that the same conclusion
about the positive definiteness of $\Ss^\nabla\omega$ on the leaves of~$\cL$
would have followed if one had merely assumed that each $u^i$
be only `strictly $\nabla$-\emph{quasi}-convex',
i.e., that $\d u^i$ be nonvanishing and $H(u^i)$ be positive definite
when restricted to the hyperplane field~$\d u^i = 0$.
Compare~\cite[Lemma~1]{MR2023129},
and the preceding discussion about their Problem~2.}
\end{proof}

This proposition provides necessary condition for the existence
of a convex Pfaff--Darboux representation.

\begin{Example}[an obstructed example]\label{ex: Obstruction}
Let $M=\mathbb{R}^n$ with standard coordinates $x=(x^i)$,
and let $\nabla$ be the (flat, torsion-free) connection
such that $\nabla(\d x^i)=0$ for $1\le i\le n$.
Let $c_i$, $f_{ij}=-f_{ji}$, and $g_{ij}=g_{ji}$
be constants and consider the $1$-form
\[
\omega = \bigl(c_i + (f_{ij}+g_{ij})x^j)\,\d x^i.
\]
Then $\d\omega= f_{ij}\,\d x^j\w\d x^i = - f_{ij}\,\d x^i\w\d x^j$
and $\Ss^\nabla\omega = g_{ij}\,\d x^i\circ\d x^j$.

Now assume that the skew-symmetric matrix $f = (f_{ij})$
has rank $2(k{-}1)<n$, so that $(\d\omega)^{k-1}\not=0$
but $(\d\omega)^{k}=0$.
Then, for generic choice of the constants $c_i$,
$\omega\w(\d\omega)^{k-1}$ will be nonvanishing on an open neighborhood $U\subset\mathbb{R}^n$ of the origin $0\in\mathbb{R}^n$,
in which case, $\omega$ will satisfy the hypotheses~\eqref{eq: omegaPfaff_k}
and~\eqref{eq: omega_nondegenerate} on~$U$.

If the symmetric matrix $g = (g_{ij})$
does not have at least $n{-}k$ positive eigenvalues,
then~$\Ss^\nabla\omega$ cannot be positive definite on
any $(n{-}k)$-dimensional subbundle $N\subset TU$, and, hence,
by Proposition~\ref{prop: ness cond}, $\omega$ cannot have a convex Pfaff--Darboux
representation in any open subset of $U$.
\end{Example}

It turns out that this necessary condition
for a local `convex' Pfaff--Darboux representation compatible
with an $\omega$-Legendrian foliation~$\cL$ is also sufficient.

\begin{Theorem}\label{thm: main}
Suppose $\nabla$ be a torsion-free affine
connection on~$M$, that $\omega\in\Omega^1(M)$
satisfy~\eqref{eq: omegaPfaff_k} and~\eqref{eq: omega_nondegenerate}
for some~$k>0$, and that $\cL$ be an $\omega$-Legendrian
foliation on~$M$ with the property that $\Ss^\nabla\omega$
pulls back to each leaf of~$\cL$ to be positive definite.
Then each $m\in M$ has an open neighborhood~$U\subset M$
on which there exist strictly $\nabla$-convex functions~$u^1,\ldots,u^k$
that are constant on the leaves of~$\cL$ in~$U$
and positive functions~$a_1,\ldots, a_k$ such that
\begin{equation*}
U^*\omega = a_1\,\d u^1 + \cdots + a_k\,\d u^k.
\end{equation*}
\end{Theorem}



% Original source lines 592--746
\begin{proof}[Proof of Theorem~\ref{thm: main}.]
There exists an $m$-neighborhood $V_0\subset M$
on which there exist smooth functions~$y^1,\ldots,y^k$
vanishing at~$m$
so that the leaves of $\d y^1 =\cdots =\d y^k = 0$
are intersections of the leaves of~$\cL$ with~$V_0$
as well as functions~$p_2,\ldots,p_k$, also vanishing
at~$m$, and a nonvanishing function~$a$ so that
\[
V_0^*\omega
= a\bigl(\d y^1 + p_2\,\d y^2 +\cdots + p_k\,\d y^k\bigr)
\]
By reversing the signs of~$a$ and the~$y^i$,
if necessary, one can assume that~$a(m)>0$.
Let~$W\subset T_mM$ be the tangent to the leaf of~$\cL$
that passes through~$m$, so that $W$ is the common kernel
of the~$\d y^i$ evaluated at~$m$.

Set~$\bar\omega = a^{-1}\omega$ and note that,
since~$\d\bar\omega \equiv a^{-1}\,\d\omega \bmod \omega$,
it follows that $\cL$ is also $\bar\omega$-Legendrian.
Moreover, since
\[
\Ss^\nabla\bar\omega = \d\big(a^{-1}\big)\circ\omega + a^{-1} \Ss^\nabla\omega,
\]
it follows that the tangent spaces of~$\cL$
(which, of course, satisfy $\omega=0$)
are also $\nabla$-positive for~$\bar\omega$.
Since~$\omega = a \bar\omega$ and $a>0$,
finding the desired convex representation for $\bar\omega$
will also yield one for~$\omega$.
Thus, it suffices to prove the theorem with~$\bar\omega$
in the place of~$\omega$, i.e., to assume that~$a=1$,
so I will do that from now on. Thus,
\[
V_0^*\omega = \d y^1 + p_2\,\d y^2 +\cdots + p_k\,\d y^k.
\]
Since~$\omega\w(\d\omega)^{k-1}\not=0$, the functions
$y^1,\ldots,y^k$, $p_2,\ldots,p_k$ have linearly independent differentials
at~$m$.

Restricting to~$V_0$, i.e., setting $M=V_0$, one has
\[
\Ss^\nabla\omega = H\big(y^1\big) + \d p_2\circ \d y^2 + \cdots + \d p_k\circ \d y^k
 + p_2 H\big(y^2\big) + \cdots + p_k H\big(y^k\big).
\]
Since the~$p_j$ vanish at~$m$,
it follows that, when restricted to~$W\subset T_mM$,
the two quadratic forms~$H\big(y^1\big)$ and $\Ss^\nabla\omega$ are equal.
Thus~$H\big(y^1\big)$ is positive definite on~$W$,
and so there is a~constant~$c>0$
so that~\smash{$H\big(y^1\big) + c \big(\d y^2\big)^2 + \cdots + c \big(\d y^k\big)^2$}
is positive definite on~$K_m = \{v\in T_mM\ \vrule\ \d y^1(v) = 0\}$.
Writing
\[
\omega = \d\bigl(y^1 + \tfrac12c \big(y^2\big)^2 + \cdots + \tfrac12c \big(y^k\big)^2\bigr)
 + \big(p_2-c y^2\big)\,\d y^2 + \cdots + \big(p_k-c y^k\big)\,\d y^k
\]
and observing that
\begin{gather*}
H\big(y^1 + \tfrac12c \big(y^2\big)^2 + \cdots + \tfrac12c \big(y^k\big)^2\big)\\
\qquad = H\big(y^1\big) + c \big(\d y^2\big)^2 + \cdots + c \big(\d y^k\big)^2
 + c y^2 H\big(y^2\big) + \cdots + c y^k H\big(y^k\big)
\end{gather*}
shows that, setting
\[
\bar y^1 = y^1 + \tfrac12c \big(y^2\big)^2 + \cdots + \tfrac12c \big(y^k\big)^2,
\qquad
\bar y^i = y^i,
\qquad\text{and}\qquad
\bar p_i = p_i - c y^i,
\]
one has
$\omega = \d\bar y^1
 + \bar p_2\,\d \bar y^2 + \cdots + \bar p_k\,\d \bar y^k$.

Thus, one could have chosen the functions~$y^1,\ldots,y^k$, $p_2,\ldots,p_k$
with~$H\big(y^1\big)$ being positive definite on the hyperplane~$K_m$.
Assume now that this was done.

It still needs to be shown that one can choose
the functions~$y^1,\ldots,y^k$, $p_2,\ldots,p_k$
with $H\big(y^1\big)$ being positive definite on all of $T_mM$,
not just on $K_m$, which is the kernel of $\d y^1$ at~$m$.
To do this, note that, if~$\phi$ is any smooth function on a neighborhood
of the origin in~$\bbR$, then
\[
H\bigl(\phi\big(y^1\big)\bigr) = \phi'\big(y^1\big) H\big(y^1\big) + \phi''\big(y^1\big) \big(\d y^1\big)^2.
\]
Hence, by choosing a~$\phi$ with $\phi(0)=0$, $\phi'(0)=1$
and $\phi''(0)>0$ sufficiently large,
I can arrange that~$\phi\big(y^1\big)$ be strictly $\nabla$-convex at~$m$.
Since
\[
\omega = \frac{1}{\phi'(y_1)} \big(\d\bigl(\phi\big(y^1\big)\bigr)
    + \phi'(y_1)p_2\,\d y^2 +\cdots + \phi'(y_1)p_k\,\d y^k\big),
\]
one sees that the functions
$\bar y^1,\ldots,\bar y^k$, $\bar p_2,\ldots,\bar p_k$, where
\[
\bar y^1 = \phi\big(y^1\big),
\qquad\text{and}\qquad
\bar y^i = y^i,
\qquad
\bar p_i = \phi'(y_1)p_i,\qquad 2\le i\le k,
\]
(with $a = 1/\phi'(y_1) >0$), give a Pfaff--Darboux representation
for~$\omega$ that is compatible with the foliation~$\cL$ and
for which $\bar y^1$ is strictly $\nabla$-convex.\footnote{This is the same idea that Ekeland and Nirenberg \cite{MR2023129}
used in their generalization of their Theorem~1
to cover the quasi-convex case.}

Thus, one can assume henceforth that,
on an open $m$-neighborhood~$V_1\subset M$,
one has a~representation of the form
\[
\omega = \d y^1 + p_2\,\d y^2 +\cdots + p_k\,\d y^k,
\]
where the functions~$y^1,\ldots,y^k,p_2,\ldots,p_k\in C^\infty(V_1)$
all vanish at~$m$,
the equations $\d y^i=0$ define the tangents to the leaves of~$\cL$ in~$V_1$,
and~$y^1$ is strictly~$\nabla$-convex.

Under these assumptions, there is a constant~$b>0$ sufficiently large
so that $H\big(y^i+by^1\big) = H\big(y^i\big)+b H\big(y^1\big)$ is positive definite at~$m$
for~$2\le i\le k$.
Thus, writing
\[
\omega = \bigl(1-b(p_2{+}\cdots{+}p_k)\bigr)\,\d y^1 + p_2\,\d\big(y^2+b y^1\big)
   + \cdots + p_k\,\d\big(y^k+b y^1\big),
\]
it follows that I can,
after restricting to an $m$-neighborhood $V_2\subset V_1$ on which
the function $a = \bigl(1-b(p_2{+}\cdots{+}p_k)\bigr)$ is positive,
dividing by~$a>0$,
and replacing~$y^j$ by $y^j+by^1$ and~$p_j$ by~$p_j/a$ for $2\le j\le k$,
assume that I have a representation
\[
\omega = \d y^1 + p_2\,\d y^2 +\cdots + p_k\,\d y^k,
\]
in which all of the~$H\big(y^j\big)$ are positive definite at~$m$,
i.e., the $y^j$ are strictly $\nabla$-convex on some neighborhood of~$m$
and the $p_i$ all vanish at~$m$.

Finally, for $\varepsilon>0$ and sufficiently small, write
\[
\omega = \d\bigl(y^1-\varepsilon\big(y^2{+}\cdots{+}y^k\big)\bigr)
   + (p_2 + \varepsilon)\,\d y^2
   + \cdots + (p_k + \varepsilon)\,\d y^k.
\]
Then, setting~$u^1 = y^1-\varepsilon\big(y^2+\cdots+y^k\big)$ and~$u^j = y^j$
for~$j>1$ and setting~$a_1 = 1$ and~${a_j = \varepsilon + p_j}$ for~$j>1$,
one achieves the desired convex Pfaff--Darboux representation
on an open $m$-neigh\-bor\-hood~$U\subset V_2$.
\end{proof}
\begin{thebibliography}{99}
\bibitem[2]{MR1083148}
Bryant R.L., Chern S.S., Gardner R.B., Goldschmidt H.L., Griffiths P.A.,
 Exterior differential systems, \textit{Math. Sci. Res. Inst. Publ.}, Vol.~18,
 \href{https://doi.org/10.1007/978-1-4613-9714-4}{Springer}, New York, 1991.

\bibitem[3]{MR1655519}
Chiappori P.-A., Ekeland I., A convex {D}arboux theorem, \textit{Ann. Scuola
 Norm. Sup. Pisa Cl. Sci.~4} \textbf{25} (1997), 287--297.

\bibitem[5]{MR2023129}
Ekeland I., Nirenberg L., A convex {D}arboux theorem, \href{https://doi.org/10.4310/MAA.2002.v9.n3.a3}{\textit{Methods Appl. Anal.}} \textbf{9} (2002), 329--344.

\bibitem[6]{MR1652709}
Zakalyukin V.M., Concave {D}arboux theorem, \href{https://doi.org/10.1016/S0764-4442(99)80092-5}{\textit{C.~R.~Acad. Sci. Paris
 S\'{e}r.~{I} Math.}} \textbf{327} (1998), 633--638.

\end{thebibliography}
\end{document}
